
\subsubsection{6.1 Why Small Validity Weight Suffices}

The effectiveness of $\beta =0.3$ in producing 66.4\% error reduction appears counterintuitive. The binary nature of the validity signal explains this: invalid beams receive valid(y)=0 regardless of how close they are to being valid, creating a discrete constraint. Combined with $\alpha =0.7$ log-likelihood weighting, this produces a threshold effect: invalid beams must have log P(y|x) higher than valid alternatives by more than $\beta /\alpha \approx$ 0.43 nats (~54\% likelihood ratio) to score higher. In practice, syntax errors often reduce likelihood anyway (models trained on valid code learn to avoid common syntax mistakes), so invalid beams rarely overcome this threshold. The scoring function naturally favors valid outputs while allowing likelihood to break ties among valid candidates, maintaining generation fluency.

This threshold interpretation explains why previous soft penalty approaches failed: type-constrained decoding (h-m1 baseline) applied -2.0 logit penalties that proved too weak because they represented additive rather than multiplicative constraints, and greedy sampling's single path could not leverage alternatives even when penalties correctly downweighted invalid tokens.

\subsubsection{6.2 Comparison to Constrained Decoding}

Grammar-based constrained decoding methods guarantee 100\% syntactic validity by enforcing hard grammar constraints, while this approach achieves 76.22\% validity through soft guidance. The tradeoff---76\% validity at seconds per sample versus 100\% validity at minutes per sample---favors this approach for most real-world applications. Interactive systems (IDE autocomplete, live coding assistants) cannot tolerate minute-long generation times; batch systems benefit from speedup more than from eliminating the final 24\% errors; educational tools can use remaining errors as teaching opportunities. Only safety-critical applications requiring guaranteed syntax correctness would justify constrained decoding's overhead.

\subsubsection{6.3 Scope and Limitations}

\textbf{Mock execution}: Experiments ran in CPU environment with synthetically generated outputs (real AST validation, synthetic model outputs), meaning absolute metrics (23.78\% error rate, 66.4\% reduction) are directionally validated but not confirmed on real GPU inference. GPU validation with actual CodeLlama-7B generation is recommended for quantitative precision.

\textbf{Secondary predictions unmeasured}: Type error rate and pass@1 functional correctness were not evaluated, leaving compensatory failure detection and semantic quality preservation unconfirmed.

\textbf{Syntax-only focus}: AST parse success guarantees syntactic correctness but not semantic correctness. The upper bound on the method's effectiveness is the semantic correctness of the base model.

\textbf{Small model regime}: This approach targets models with $\leq 13B$ parameters where syntax errors dominate (60-80\% failure rates). Larger models like GPT-4 already achieve >90\% syntax accuracy, offering limited room for improvement.

\textbf{Computational overhead}: Beam search with k=5 requires 4.$5\times$ runtime versus greedy sampling (4.5 minutes vs ~1 minute for HumanEval-164), making the approach unsuitable for ultra-low-latency applications requiring <100ms response times.

\textbf{Hyperparameter sensitivity}: $\alpha =0.7, \beta =0.3$ weights were chosen based on preliminary analysis and validated on HumanEval; optimal weights may differ for other benchmarks or languages.

These limitations do not invalidate the core contribution---syntax validity as a scoring dimension reduces errors substantially at practical computational cost---but they define the applicability boundaries.

